
%%%%%%%%%%%%%%%%%%%%%%%%
%
% $Autor: Gerken $
% $Pfad: GDV/Vortraege/latex - Ausarbeitung/Kapitel/Einleitung.tex $
% $Version: 1 $
%
%%%%%%%%%%%%%%%%%%%%%%%%

\chapter{\textcolor{red}{Introduction}
	
	Eine WordClock (zu Deutsch: Wortuhr) bricht mit der traditionellen Darstellung der Zeit durch Zeiger oder Ziffern. Stattdessen nutzt sie eine Matrix aus Buchstaben, die gezielt hinterleuchtet werden, um die Uhrzeit in Form von geschriebenen Sätzen anzuzeigen - zum Beispiel "Es ist viertel nach zehn" oder "Es ist halb eins". Diese Art der Anzeige verwandelt einen funktionalen Zeitmesser in ein modernes Designobjekt, das Information und Ästhetik auf besondere Weise verbindet.
